\SourceLeaf{308}
\fancyhead[L]{\small UNTERSUCHUNGEN ÜBER THEMEN DER HÖHEREN GEODÄSIE}
\renewcommand{\Article}[1]{\par\medskip\Needspace{4\baselineskip}\pdfbookmark[1]{Art. #1}{geodaesie2-art-#1}\centerline{\textbf{#1.}}\nobreak\smallskip}
\begin{center}
UNTERSUCHUNGEN\par\bigskip
ÜBER\par\bigskip
{\large THEMEN DER HÖHEREN GEODÄSIE.}\par\medskip
\rule{.65in}{.3pt}
\end{center}
\noindent\hspace*{\parindent}Zu den Hauptaufgaben der höheren Geodäsie gehört es, aus der Länge einer Dreiecksseite auf der Erdoberfläche, dem Azimut an einem ihrer Endpunkte und der geografischen Breite dieses Endpunkts das Azimut am anderen Endpunkt, dessen Breite und den Längenunterschied beider Punkte abzuleiten. Für die Kugeloberfläche ist der Zusammenhang zwischen diesen sechs Größen am Schluss der ersten Abhandlung in der einfachsten und für die genaueste Rechnung geeigneten Form angegeben. Diese Form lässt sich auch leicht zur bequemen Lösung der Aufgabe selbst verwenden. Daraus ergibt sich der Wunsch nach einer entsprechenden Lösungsart, die unmittelbar für die Ellipsoidoberfläche gilt. Eine solche zu entwickeln ist das Ziel der vorliegenden Abhandlung. Zuvor soll jedoch die Lösung für die Kugeloberfläche noch deutlicher dargestellt werden. Um Verweise zu erleichtern, setze ich die Nummerierung der Artikel aus der ersten Abhandlung fort.

\Article{18}
Um den Grad der Genauigkeit der Formeln aus Artikel 17 besser beurteilen zu können, muss man noch die Glieder der nächsthöheren Ordnung entwickeln. Es lohnt sich jedoch,
